\documentclass[10pt,twocolumn]{article}
\usepackage[T1]{fontenc}
\usepackage[utf8]{inputenc}
\usepackage{lmodern}
\usepackage[margin=1.8cm,top=2.4cm,bottom=2cm,columnsep=0.6cm]{geometry}
\usepackage{fancyhdr}
\usepackage{titlesec}
\usepackage{booktabs}
\usepackage{amsmath,amssymb,amsfonts}
\usepackage{microtype}
\usepackage{xcolor}
\usepackage{mdframed}
\usepackage{parskip}
\usepackage{enumitem}
\usepackage{hyperref}

\definecolor{rulered}{RGB}{180,30,30}
\definecolor{boxgray}{RGB}{245,245,245}
\definecolor{keywordgray}{RGB}{100,100,100}

\pagestyle{fancy}
\fancyhf{}
\fancyhead[L]{\textsc{\small Claude Mag}}
\fancyhead[C]{\small\textit{Machine Learning Research Digest}}
\fancyhead[R]{\small 2026-07-15}
\fancyfoot[C]{\small\thepage}
\renewcommand{\headrulewidth}{0.8pt}

\titleformat{\section}{\normalsize\bfseries\color{rulered}}{}{0pt}{}%
  [\vspace{-4pt}\color{rulered}\rule{\columnwidth}{0.4pt}]
\titlespacing{\section}{0pt}{8pt}{4pt}

\hypersetup{colorlinks=true,urlcolor=rulered,linkcolor=black}

\begin{document}
\setlength{\parindent}{0pt}
\setlength{\parskip}{4pt}

{\small\color{keywordgray}\textbf{Keywords:}
  mechanistic interpretability \textbullet{}
  coding agents \textbullet{}
  probing \textbullet{}
  latent representations}\par\vspace{4pt}

{\LARGE\bfseries\raggedright
  Latent Programming Horizons}\par\vspace{4pt}

{\large\itshape\raggedright
  Coding Agents Know the Outcome of Their Edits
  Up to 25 Steps Before They Make Them}\par\vspace{6pt}

{\small\textbf{By} Andr\'e Silva, Han Tu \& Martin Monperrus
  (KTH Royal Institute of Technology, Stockholm)}\par\vspace{2pt}

{\small\color{keywordgray}arXiv:2607.05188 \textbullet{}
  July 6, 2026}
\par\vspace{2pt}\noindent{\color{rulered}\rule{\columnwidth}{0.6pt}}\vspace{6pt}

When a language model acts as a coding agent, it spends dozens of
steps reading files, reasoning about bugs, writing patches, and
running tests---a process that is opaque from the outside. This paper
opens that black box. Using logistic-regression probes on the residual
streams of GPT-4o and Claude 3.7 Sonnet, the authors show that the
model's internal representations linearly encode the correctness of
the program being edited with AUC up to 0.83---and that these
representations run up to 25 steps \emph{ahead} of the agent's
own actions, constituting what the authors call a \textbf{latent
programming horizon}.

\begin{mdframed}[backgroundcolor=boxgray,linecolor=rulered,linewidth=1pt,
    innerleftmargin=6pt,innerrightmargin=6pt,
    innertopmargin=6pt,innerbottommargin=6pt]
\textbf{\small AT A GLANCE}\par\vspace{4pt}
\begin{description}[leftmargin=0pt,labelsep=4pt,itemsep=2pt,parsep=0pt,
    font=\small\bfseries]
  \item[Problem:] The internal states of coding agents during multi-step
    repair episodes are not understood; it is unknown whether models
    plan ahead or react greedily.
  \item[Why it matters:] A latent horizon implies agents have exploitable
    internal signals for early stopping, intervention, and lookahead
    guidance without external execution.
  \item[Method:] Logistic-regression probes on final-layer hidden states
    recorded at each step of coding-agent episodes; future probes
    predict outcomes $k$ steps ahead.
  \item[Result:] AUC 0.83 for correctness at step $t$; horizon $k^*$
    up to 25 steps on SWE-bench Verified.
  \item[Evidence:] Two models (GPT-4o, Claude 3.7), two benchmarks
    (SWE-bench Verified, HumanEval-Pro); cross-benchmark transfer
    retains AUC $>$ 0.70.
\end{description}
\end{mdframed}

\section{The Big Picture}

Large language models acting as software engineering agents can resolve
real GitHub issues, pass coding interviews, and complete complex
multi-file refactors---but the mechanism by which they succeed or fail
across a long episode remains opaque. Mechanistic interpretability
research has largely focused on single forward passes or small
fixed-context tasks.

Coding-agent episodes are different: the model sees an evolving
repository state, modifies files, receives test outputs, and makes
new decisions over potentially hundreds of steps. Understanding what
the model ``knows'' at each step has direct practical value---if the
model's internal state already encodes the likely outcome of the episode,
that signal could be used to guide, interrupt, or improve the agent
without waiting for external test execution.

This paper introduces the first systematic mechanistic interpretability
study of coding agents over full multi-step episodes.

\section{Why Existing Approaches Fall Short}

Prior mechanistic interpretability work probes fixed-context forward
passes (e.g., factual recall, in-context learning, arithmetic). These
methods cannot be applied directly to coding-agent episodes because
(1)~the context evolves across steps as files are modified, (2)~ground
truth labels depend on external execution of the evolving codebase, and
(3)~the episode length (up to 200 steps) is far longer than tasks
previously studied.

Behavioral analysis of coding agents (e.g., measuring pass@k or
success rate) describes what agents do but not what they represent.
Existing evaluation frameworks do not distinguish between agents that
reason correctly but act incorrectly and agents that reason incorrectly
and happen to succeed by luck---a distinction this paper makes tractable
via probing.

\section{Core Mechanism}

The latent programming horizon arises because the model processes
the full current repository state and reasoning history at each step,
and its internal representations must encode relevant aspects of the
program's state to generate useful next-action tokens.

The key empirical finding is that correctness-relevant features are
not just encoded in the final step before an edit---they are
stably encoded up to 25 steps earlier in the episode. This suggests
the model is tracking a ``mental model'' of the program's correctness
that is both accurate and temporally extended.

\section{Technical Formulation}

At each step $t$ of an episode, the last-position final-layer residual
stream vector $h_t \in \mathbb{R}^{4096}$ is recorded. Binary labels
$y_t \in \{0, 1\}$ are obtained by running pytest on the repository
state after each agent edit (parsed: 1/0; test-pass: 1/0; etc.).

For immediate probes, a logistic regression $f_\theta(h_t) \to \hat{y}_t$
is trained on $(h_t, y_t)$ pairs.

For future probes at horizon $k$, the same probe structure predicts
$y_{t+k}$ from $h_t$. The latent programming horizon is:
\[
k^* = \max\bigl\{k\;:\;
  \mathrm{AUC}(f_{t \to t+k}) \geq 0.55,\;p < 0.05\bigr\}
\]
where significance is assessed by paired permutation test over episodes.

\section{Learning or Inference Procedure}

\begin{enumerate}[itemsep=2pt,parsep=0pt]
  \item Run coding agents on SWE-bench Verified (300 issues) and
    HumanEval-Pro (150 problems) with step-by-step residual stream
    logging enabled via API hooks.
  \item After each edit, execute sandboxed pytest and record binary
    labels for four tasks (parses, passes tests, reduces failures,
    introduces regression).
  \item Train logistic-regression probes on 80\% of episodes;
    evaluate on 20\% held-out episodes.
  \item Sweep horizon $k$ from 0 to 40; measure AUC and $p$-value
    at each horizon to determine $k^*$.
  \item Perform cross-benchmark transfer: train on SWE-bench,
    evaluate on HumanEval-Pro without retraining.
\end{enumerate}

\section{What the Guarantee Says}

The core finding is not a guarantee in the theoretical sense, but a
robust empirical regularity: across two models and two benchmarks,
the latent programming horizon $k^*$ is consistently 20--28 steps
for ``passes tests'' and 12--16 steps for ``introduces regression.''
Probes trained on one model do not transfer to the other, indicating
the representations are model-specific even if the phenomenon is
model-agnostic. Cross-benchmark transfer holds above AUC 0.70 for all
tasks, confirming the probe captures generic code-quality structure.

\section{Experimental Findings}

\begin{center}
\small
\begin{tabular}{llcc}
\toprule
Task & Model & AUC ($k{=}0$) & Horizon $k^*$ \\
\midrule
Parses & GPT-4o & 0.81 & 28 \\
Passes tests & GPT-4o & 0.78 & 22 \\
Reduces fails & Claude 3.7 & 0.83 & 25 \\
Introduces reg. & Claude 3.7 & 0.71 & 14 \\
\bottomrule
\end{tabular}
\end{center}

Cross-benchmark transfer (SWE-bench $\to$ HumanEval-Pro) retains
AUC above 0.70 for all tasks. Middle-layer probes (layer 24/48)
achieve only AUC 0.63, indicating the relevant representations are
built through the full network depth. Non-linear probes improve
AUC by only 0.02, confirming approximate linear decodability.

\section{Ablations and Interpretation}

\textbf{Layer depth:} Probing earlier layers degrades AUC monotonically
from 0.83 (final layer) to 0.51 (first layer), suggesting correctness
representations are progressively constructed across transformer depth.

\textbf{Non-linearity:} Adding one MLP hidden layer improves AUC by
0.02 on average. The dominant signal is linear, consistent with the
probing literature on factual representations.

\textbf{Shuffled-label control:} Probes trained on permuted labels
achieve AUC $0.50 \pm 0.01$, ruling out statistical artifacts from
episode length or context size.

\textbf{Practical implication:} The 25-step horizon means an
external monitor could identify likely-failing episodes 25 steps
early and redirect the agent---a potential 25-step saving per
failed episode. On SWE-bench Verified, where 40\% of episodes fail,
this would reduce total token cost by $\sim$10\%.

\section*{Reference}

Andr\'e Silva, Han Tu, Martin Monperrus.
\textbf{Latent Programming Horizons in Coding Agents}.
arXiv:2607.05188, July 2026. KTH Royal Institute of Technology.
\url{https://arxiv.org/abs/2607.05188}

\end{document}
